\section{Conclusion}
\label{sec:conclusion}

We have introduced \pt{} redistricting, a Search-layer algorithm that runs
$N_{\mathrm{replicas}}$ \fr{} chains simultaneously at different balance tolerances
arranged on a geometric ladder.
Adjacent chains periodically propose to exchange their current plans via the
Metropolis criterion; accepted swaps inject global diversity from hot chains
into the cold chain, enabling escape from local optima without modifying the
cold chain's distributional target.

\paragraph{Key design choices.}
Three design choices make \pt{} redistricting practical:
\begin{enumerate}
\item The \emph{geometric ladder} $\varepsilon_i = \varepsilon_c \times
  (\varepsilon_h/\varepsilon_c)^{i/(N-1)}$ provides approximately equal swap
  acceptance rates across all adjacent pairs, calibrated by the ${\approx}23\%$
  empirical target from continuous systems~\citep{kone2005}.
\item The \emph{swap interval} of 10 steps balances chain diffusion time
  against swap frequency, consistent with established tempering practice~\citep{earl2005}.
\item The \emph{cold chain record collection} stores all $T+1$ cold chain
  states, sorts by EC, and returns the plan at the requested percentile.
  This allows percentile-based plan selection without rerunning the chain.
\end{enumerate}

\paragraph{Empirical summary.}
PT($N=4$) achieves 7.6\% EC reduction over single-chain \fr{} on Texas ($k=38$),
3.4\% on Wisconsin ($k=8$), and 3.3\% on North Carolina ($k=14$).
The improvement is largest where local-optimum trapping is most severe.
Observed swap acceptance rates of $18$--$24\%$ confirm that the geometric
ladder is well-calibrated across all three states.

\paragraph{Position in the algorithm landscape.}
\pt{} is a Search-layer algorithm, orthogonal to the Structure layer
(\geosec{}, \ar{}, \sa{}) and to the WeightsOverride layer.
It composes with any structure/weights configuration.
The recommended decision rule for practitioners is:
\begin{itemize}
\item \textbf{$k \le 20$}: Use single-chain \fr{} or \cs{}. PT overhead not
  justified for adequate mixing states.
\item \textbf{$k \ge 30$, auditability priority}: Use PT ($N=4$ default,
  $N=6$ for TX/CA).
  The SHA-256 seed derivation provides a complete audit chain from
  \texttt{base\_seed} to the selected plan.
\item \textbf{$k \ge 30$, geographic structure strong}: Consider \ms{} or
  PT+\ms{} (future work).
\end{itemize}

\paragraph{Implementation.}
\pt{} redistricting is implemented in the \texttt{redist} binary under
\texttt{--search parallel-tempering}, with parameters \texttt{--pt-replicas},
\texttt{--pt-swap-interval}, \texttt{--pt-cold-tol}, and \texttt{--pt-hot-tol}.
The \texttt{ParallelTemperingChain} struct lives in
\texttt{redist-ensemble/src/parallel\_tempering.rs} and exposes
\texttt{swap\_acceptance\_rate()}, \texttt{cold\_chain\_ec\_mean()},
and \texttt{select\_plan(p)} as the primary API.

\paragraph{Open questions.}
\begin{enumerate}
\item \textbf{Parallel execution.}
  The $N$ chains at different tolerances are independent between swap steps---an
  embarrassingly parallel workload.
  Rayon-based parallel execution would reduce wall-clock cost from $N\times$
  to approximately $1\times$ the single-chain runtime, making PT competitive
  on wall-clock time even for large $N$.
  This is deferred as a Phase~2 implementation task.
\item \textbf{Mixed-chain design.}
  Should hot chains use standard \recom{} (faster per step, higher acceptance)
  while only the cold chain uses \fr{} (distributional correctness)?
  Mixed-chain designs complicate the swap Metropolis criterion but may improve
  wall-clock performance.
\item \textbf{Adaptive replica count.}
  Adding or removing replicas during the run based on observed swap acceptance
  rates would optimise the ladder dynamically.
  This requires careful audit-chain handling to maintain reproducibility.
\item \textbf{PT $\times$ Multi-Scale.}
  Each PT replica could run an \ms{} chain at its tolerance.
  This would combine the global diversity injection of PT with the large-proposal
  benefit of \ms{}, potentially providing the best of both mechanisms for
  TX-class states.
\item \textbf{Non-adjacent swaps.}
  Current implementation only proposes swaps between adjacent pairs $(i, i+1)$.
  Non-adjacent swaps (between any pair of replicas) may improve global mixing
  but require adjusting the Metropolis criterion and are less standard in the
  tempering literature.
\end{enumerate}

These extensions inherit the clean three-layer compositor design of the existing
platform and require no structural changes to the bisection tree or the
\fr{} chain implementation.
